%% kapa: exam=1 round=31 year=2020 period=1 form=A subject=civil_law subject_ko=민법 seq=35 number=35
%% source: Q-Net 감정평가사 기출문제 (cst00309) artlSeq=5207219 file=extracted/1008/5207219/A형/A형/2020년도 제31회 감정평가사 1차 시험_1교시_A형.pdf
%% conversion: kordoc 4.24.51 -> tools/gen1.py (auto); review=proofread-sonnet; flags=none
\begin{question}{35}
지역권에 관한 설명으로 옳은 것은? (다툼이 있으면 판례에 따름)
\begin{choices}
  \item 지역권은 점유를 요건으로 하는 물권이다.
  \item 지역권은 독립하여 양도ㆍ처분할 수 있는 물권이다.
  \item 통행지역권은 지료의 약정을 성립요건으로 한다.
  \item 통행지역권의 시효취득을 위하여 지역권이 계속되고 표현되면 충분하고 승역지 위에 통로를 개설할 필요는 없다.
  \item 통행지역권을 시효취득한 요역지소유자는, 특별한 사정이 없으면 승역지의 사용으로 그 소유자가 입은 손해를 보상하여야 한다.
\end{choices}
\end{question}
\answer{5}
